\scanpage{31}
\mspage{30}

% L1–5
\deleted{Ainsi} Soit $\deleted{U}X=\deleted{U}X_{\Phi}$ une partie modérée. Un $x\in L$ est dit en pos. rel. modérée avec $\deleted{U}X$, s'il l'est avec les $T_i$ (\inserted{\uline{tronçons}} \uline{composants} de $\deleted{U}X$) i.e. s'il l'est avec $\underset{\partial U}{\underline{\bigcup\partial T_i}}$, ou encore avec $\operatorname{Fr}(\deleted{U}X)$.
% Lleft margin, alongside lines 1–13
% slanted
\margin{Plus \unsure{généralemt|généralement}, Si $\deleted{U}X\unsure{,|\cap}\deleted{\unsure{V|U'}}Y$ sont modérées, on dit que $\deleted{U}X,Y$ sont en pos. rel. modérée si \deleted{$\partial U$ et $U'$}, $\partial\deleted{U}X,\partial\deleted{V}Y$ le sont ou encore $\operatorname{Fr}\deleted{U}X,\operatorname{Fr}Y$ le sont. Ex: si $\deleted{U}X\subset Y$, alors $(\deleted{U}X,\deleted{V}Y)$ \deleted{\unsure{est}} \unsure{modérée|modéré}.}
% L5–7
Par exemple \deleted{Soit} $x\in\deleted{U}X$, ou $x\in\deleted{\unsure{\operatorname{Fr}|\partial}}\partial U$, à fortiori de $\operatorname{Fr}(\deleted{U}X)$, a cette propriété. Soit
% L8–9
\[
\begin{aligned}
\deleted{U}X'=C_{\mathrm{mod}}\unsure{U U)|(U)|U)}=\bigl\{x\in L\setminus U\mid{}&x\text{ est en pos. rel. mod.}\\
&\text{p.r. à }U\bigr\}
\end{aligned}
\]
% L10–13
Alors \deleted{\unsure{$\mathcal E$|$\mathcal C$|$C$}} $\deleted{U}X'$ est \unsure{modérée|modéré}, et c'est la plus grande partie modérée $\deleted{V}Y$ de $L$ telle que $\deleted{V}Y\subset L\setminus\deleted{U}X$, et $Y$ en position relative modérée par. à $X$. Si
% L14
\[
\{T_1\ll T_2<\ldots\ll T_n\}
\]
% L15–20
est le drapeau définissant \deleted{\unsure{$T$|$\Phi$}} $X$, $X'$ est défini par le drapeau formé des $n-1$ tronçons interstitiels \deleted{entre} pour $(T_1,T_2)$, $(T_2,T_3)$ \ldots\ $(T_{n-1},T_n)$, plus (dans la mesure où ils existent) les tronçons, interstitiel gauche de $T_1$, et le tronçon interstitiel droit du $n$. \unsure{Suite|D'autre}
